\documentclass[10pt,a4paper]{amsart}
\usepackage[margin=19mm]{geometry}
\usepackage{amsmath,amssymb,mathtools}
\usepackage[scr=rsfs,cal=euler]{mathalfa}
\usepackage{xfrac}
\usepackage[unicode,hidelinks,bookmarks=false]{hyperref}
\newcommand{\ie}{i.e.\ }
\newcommand{\cf}{cf.\ }
\newcommand{\resp}{respectively\ }
\newcommand{\Subsection}{\subsection}
\providecommand{\nobreakdash}{\nobreak\hbox{-}}
\setlength{\emergencystretch}{2em}
\multlinegap0pt
\usepackage{tensor}
\usepackage{amssymb}
\usepackage{enumerate}
\usepackage{tikz}
\newtheorem{prop}{Proposition}
\def\H#1{\mathcal{H}(#1)}
\def\N{\mathbb{N}}
\def\weakr{\leq_{R}}
\title{On the Weak Right Order of Right-Angled Coxeter Systems}
\author{Harrison Gimenez}
\date{Source: arXiv:2606.21104v1 (CC BY 4.0)}
\begin{document}
\maketitle

\noindent\emph{Source and licence.} On the Weak Right Order of Right-Angled Coxeter Systems. Version arXiv:2606.21104v1 (CC BY 4.0), distributed by arXiv under the Creative Commons Attribution 4.0 International licence, http://creativecommons.org/licenses/by/4.0/, as recorded in the arXiv licence metadata for this exact version and shown on the abstract page https://arxiv.org/abs/2606.21104v1. Full licence notice: \texttt{licenses/arxiv-2606.21104-CC-BY-4.0.txt}.\par\smallskip

\noindent\textit{Editorial scope.} Counted unit: the article's prop joinisinvolution (source0.tex lines 489--491). The article's complete original proof of it is delivered with it (source0.tex lines 493--517, one \texttt{proof} environment of 25 lines). No mathematics is written, repaired or re-derived: the statement and the proof are the article's own lines, in the article's own notation, and no retained line is substituted or re-flowed.  Retained uncounted as the source-local context the proof needs: lines 338--353; lines 371--385; lines 393--401; lines 411--413; lines 424--432; lines 472--481.  Nothing was omitted from any retained span, so the package carries no omission record.  No retained line contains a citation, so the wrapper carries no bibliography and no bibliography entry.  No retained line contains a figure environment, an image-inclusion command or drawing code, so no image is delivered and the package declares no asset dependency.  Editorial formatting only, in addition: an AMS \texttt{amsart} preamble that restates the article's own declarations and macros used by the retained lines, namely the environments \texttt{prop} on separate counters, together with the standard packages those lines need.  The retained environments print in this document as Proposition 1 in order of appearance, whereas the article prints the same items with the numbering of its own full document.  The complete native source of the article as distributed in the arXiv e-print for this exact version is bundled unchanged as \texttt{sources/203256/source0.tex}.
\begin{prop} \label{removemaxprop}
    Let $ (W,S)$ be a Coxeter system. Let $ w$ be a fully commutative element. Consider $\H{w}$. Then the following are true:

    \begin{enumerate}
        \item Let $ D_{R}(w) : = \{ s\in S \mid \ell(ws) < \ell(w)  \}$.
        Then $ s\in D_{R}(w)$ if and only if there is a maximal vertex $ q\in \H{w}$ such that the $S$-label of $ q$ is $s$.

        \item Let $ s\in D_{R}(w)$. Then $ ws$ is fully commutative, and $ \H{ws}\cong \H{w} \setminus \{ q \}$ for some maximal vertex $q\in \H{w}$ such that the $S$-label of $q$ is $s$.

        \item Let $ q_{1}, q_{2} \in \H{w}$ denote two distinct maximal vertices. Then the $S$-labels of $ q_{1}$ and $q_{2}$ must be distinct.

        \item Let $ q \in \H{w}$ be a maximal element with $ S$-label $s$. Then $ ws$ is fully commutative and $ \H{ws} \cong \H{w} \setminus \{ q  \}$.

        \item Let $ x,y\in W$ such that $ x$ and $ y$ are both fully commutative elements. Then $ x=y$ if and only if $ \H{x} \cong \H{y}$.
    \end{enumerate}
\end{prop}

\begin{prop} \label{removeminprop}
    Let $ (W,S)$ be a Coxeter system. Let $ w$ be a fully commutative element. Consider $\H{w}$. Then the following are true:

    \begin{enumerate}
        \item Let $ D_{L}(w) : = \{ s\in S \mid \ell(sw) < \ell(w)  \}$.
        Then $ s\in D_{L}(w)$ if and only if there is a minimal vertex $ p\in \H{w}$ such that the $S$-label of $ p$ is $s$.

        \item Let $ s\in D_{L}(w)$. Then $ sw$ is fully commutative, and $ \H{sw}\cong \H{w} \setminus \{ p \}$ for some minimal vertex $p\in \H{w}$ such that the $S$-label of $p$ is $s$.

        \item Let $ p_{1}, p_{2} \in \H{w}$ denote two distinct minimal vertices. Then the $S$-labels of $ p_{1}$ and $p_{2}$ must be distinct.

        \item Let $ p \in \H{w}$ be a minimal element with $ S$-label $s$. Then $ sw$ is fully commutative and $ \H{sw} \cong \H{w} \setminus \{ p  \}$.

    \end{enumerate}
\end{prop}

\begin{prop} \label{existscorrprop}
    Let $ (W,S )$ be a Coxeter system, and let $ w\in W$ be a fully commutative element. Let $I \subseteq \H{w}$ be a labeled order ideal of $ \H{w}$. Then there exists a unique $ x\in W$ such that the following are true:

    \begin{enumerate}
        \item $ x \weakr w$, and $ x$ is fully commutative

        \item $ \H{x} \cong I$
    \end{enumerate}
\end{prop}

\begin{prop} \label{mustbecompprop}
    Let $ (W,S)$ be a Coxeter system, and let $ w\in W$ be fully commutative. Let $ q_{1}, q_{2} \in \H{w}$ be two distinct elements with the same $S$-label. Then $ q_{1} \leq_{w} q_{2}$ or $ q_{2} \leq_{w} q_{1}$.
\end{prop}

\begin{prop} \label{equividealprop}
    Let $ (W,S)$ be a Coxeter system. Let $ w\in W$ be a fully commutative element. Let $x\in W $. Then the following two conditions are equivalent:

    \begin{enumerate}
        \item $ x \weakr w$

        \item $ x$ is fully commutative, and there exists a labeled order ideal $ I \subseteq \H{w}$ such that $ \H{x} \cong I$ (as $S$-heaps modulo isomorphism).
    \end{enumerate}
\end{prop}

\begin{prop} \label{joinidealprop}
    Let $ (W,S)$ be a Coxeter system. Let $ w\in W$ be fully commutative. Suppose that there exist fully commutative elements $ x,y \in W$ and labeled order ideals $ I , J \subseteq \H{w}$ such that

    \begin{itemize}
        \item $ \H{x} \cong I$ and $ \H{y} \cong J$

        \item $ \H{w} = I \cup J$
    \end{itemize}
    Then $ w = x \vee y$ where $ \vee $ denotes the join in the weak right order of $(W,S)$.
\end{prop}

\begin{prop} \label{joinisinvolution}
    Let $ (W,S)$ be a Coxeter system, and let $ w\in W$ be fully commutative. Let $ x,y\in W$ be such that $ x \weakr w$ and $ y\weakr w$. Suppose further that $ x^{2} = 1$ and $ y^{2}=1$. Then $ x\vee y$ exists and $ (x \vee y)^{2}  =1$.
\end{prop}

\begin{proof}
    Since $ x \weakr w$ and $ y \weakr w$, there exists labeled order ideals $I, J \subseteq \H{w}$ such that $\H{x} \cong I $ and $ \H{y} \cong J$ by Proposition \ref{equividealprop}. The union of two labeled order ideals is also a labeled order ideal, so $I \cup J \subseteq \H{w}$ is a labeled order ideal. By Proposition \ref{existscorrprop}, there exists some $ z \weakr w$ such that $ \H{z} \cong I \cup J$. By Proposition \ref{joinidealprop}, we must have that $ z = x \vee y$. Since $ \H{x\vee y} \cong I \cup J$, let us naturally view $ I$ and $J$ as labeled order ideals of $ \H{x\vee y}$ and thus $ \H{x \vee y} = I \cup J$.

    We now prove the claim that $ (x\vee y)^{2} = 1$ via induction on $n= \ell(x) + \ell(y)$. Note that if $ n = 0$, then $ x=1$ and $ y=1$, and hence $ x\vee y = 1$, and thus the desired conclusion is true.

    Suppose now that the claim holds for all $ k =0,1,2,  \dots , n$ where $ k = \ell(x) + \ell(y)$. We wish to prove that the claim holds for $ n+1 = \ell(x) + \ell(y)$. Note that because $ \ell(x) + \ell(y) = n+1$ where $n\in \N$, it follows that at least one of $ x$ or $y$ is not the identity, and hence $ x\vee y$ is not the identity. Thus, $ \H{x \vee y} \neq \emptyset$. 
    
    Let us suppose that all vertices of $ \H{x \vee y}$ are minimal. By Proposition \ref{removeminprop} part (3), this would imply that each vertex has a distinct $S$-labeling. Furthermore, by definition of $ \leq_{x\vee y}$, if $ r$ and $s$ were the $S$-labels of distinct vertices in $ \H{x\vee y}$, then because all vertices are minimal, we would necessarily have that $ m(r,s)  =2$ (if $m(r,s)\neq 2$, then this would imply that one vertex is strictly greater than the other with respect to $ \leq_{x\vee y}$, contradicting minimality of all vertices). But because $ m(r,s) = 2$, this implies that $ (rs)^{2} = 1$ and hence $ rs=sr$. Thus, if all vertices of $ \H{x\vee y}$ are minimal, then $ x\vee y$ would be a product of distinct simple reflections that pairwise commute. Therefore, we would have that $ (x\vee y)^{2} = 1$. Thus, from now on, we can assume that not all vertices of $ \H{x\vee y}$ are minimal. In particular, since we are dealing with finite posets, there must exist some maximal vertex $ q \in \H{x\vee y}$ such that $ q$ is not minimal in $ \H{x\vee y}$. Since $ \H{x\vee y} = I \cup J$, we must have that $ q\in I$ or $ q\in J$. Without loss of generality, let us assume that $ q\in I$. From now on, let $ s\in S$ denote the $S$-labeling of $q$. Since $ q$ is maximal in $ \H{x\vee y} = I \cup J$ where $ q\in I$, it follows that $ q$ is also maximal in $I$. But if $ q$ is maximal in $I$ and $ \H{x} \cong I$, then it follows that $ s\in D_{R}(x)$ by Proposition \ref{removemaxprop} part (1). But since $ x^{2}=1$, we have that $ D_{R}(x) = D_{L}(x)$, so we conclude that there must exist some minimal vertex $ p \in I$ such that the $S$-label of $p$ is $s$. Because $ I$ is an order ideal of $\H{x\vee y}$, it follows that $p$ is also minimal in $ \H{x\vee y}$. Note that $ p \neq q$ since $ q$ is assumed to be maximal but not minimal in $ \H{x\vee y}$. Furthermore, since $p$ and $ q$ have the same $S$-label in $ \H{x\vee y}$ it follows that $p <_{x\vee y} q$ by Proposition \ref{mustbecompprop}.

    \textbf{Case 1: $p\notin J$}
    
    
    Suppose that $ p \notin J$. Not that since $J$ is a labeled order ideal, and because $ p <_{x\vee y} q$, it follows that $ q\notin J$. Hence, $ \H{x\vee y} \setminus \{ p,q \} = (I \setminus \{  p,q \} ) \cup J$. By applying Proposition \ref{removemaxprop} part (4) and then applying Proposition \ref{removeminprop} part (4), we deduce that $ \H{s(x\vee y)s} \cong \H{x\vee y} \setminus \{ p,q \} $. Note also that $ (I \setminus \{  p,q \} )$ and $J$ are still labeled order ideals of $\H{s(x\vee y)s} \cong \H{x\vee y} \setminus \{ p,q \} $. Furthermore, by applying Proposition \ref{removemaxprop} part (4) and then applying Proposition \ref{removeminprop} part (4), we deduce that $ \H{sxs} \cong I \setminus \{ p,q \}$. We still have $ \H{y} \cong J$. Thus, since $ \H{s(x\vee y )s} = (I \setminus \{  p,q \} ) \cup J$, we can apply proposition \ref{joinidealprop} to conclude that $ s(x\vee y)s = (sxs)\vee y$. But note that $ s(x\vee y)s$ is a fully commutative element where $ sxs \weakr s(x\vee y)s$, $ y \weakr s(x\vee y)s$, $ (sxs)^{2}=1$, $ y^{2}=1$, and $ \ell(sxs) + \ell(y) = \ell(x) -2 +\ell(y) =n+1 - 2 = n-1< n+1$. Hence, we can apply induction to conclude that $ (s(x\vee y)s)^{2} = 1$, which implies $ (x\vee y)^{2} = 1$. 

    \textbf{Case 2: $ p \in J$ and $ p$ is maximal in $J$}

    Suppose now that $ p\in J$ and $ p$ is maximal in $J$. Since $p$ is maximal in $J$, $ q\notin J$. Note also that because $ p $ is minimal in $\H{x\vee y}$, it follows that $ p$ is minimal in $J$. But because $p$ is both maximal and minimal in $J$, it follows that $p$ is incomparable to all other vertices of $ J$. Let $ r$ denote the $S$-label of some arbitrary vertex in $ J \setminus \{ p \}$. Note that because $ p$ is incomparable to all vertices in $ J \setminus \{ p \}$, it follows that $ m(r,s) = 2$ (for if $ m(r,s) \neq 2$ for some $ r$ where $r$ is an $S$-label of some vertex in $ J \setminus \{ p \}$, then $ p$ would be comparable to some element in $ J \setminus \{ p \}$ by definition of $ \leq_{x\vee y}$). Since $ \H{y} \cong J$, this implies that in any given reduced expression for $y$, the simple reflection $ s$ will commute with any other simple reflection appearing in that reduced expression. In particular, $ sy=ys$. Note that since $ p\in J$ and $ q\notin J$, we have that $ \H{x\vee y} \setminus \{  p,q \} = (I \setminus \{ p,q \}) \cup (J \setminus \{ p\})$. Note that $I \setminus \{  p,q\} $ and $ J \setminus \{ p\}$ are labeled order ideals of $ \H{x\vee y} \setminus \{  p,q \}$. By applying Proposition \ref{removemaxprop} part (4) and then applying Proposition \ref{removeminprop} part (4), we deduce that $ \H{s(x\vee y)s} \cong \H{x\vee y} \setminus \{ p,q \} $. Similarly, $ \H{sxs} \cong I \setminus \{ p,q \}$. By applying Proposition \ref{removemaxprop} part (4), we also deduce that $ \H{ys} \cong J \setminus \{ p \}$. Hence, we can apply Proposition \ref{joinidealprop} to conclude that $ s(x\vee y)s = (sxs) \vee (ys)$. But note that $ s(x\vee y)s$ is a fully commutative element where $sxs \weakr s(x\vee y)s $, $ ys \weakr s(x\vee y)s$, $ (sxs)^{2}  =1$, $ (ys)^{2} = 1$, $ \ell(sxs)+ \ell(ys) = \ell(x) - 2 + \ell(y) -1 = n+1 -3 = n-2< n+1$. Hence, we can apply induction to conclude that $ (s(x\vee y)s)^{2} = 1$, which implies that $ (x\vee y)^{2} = 1$.

    \textbf{Case 3: $p \in J$ and $ p$ is not maximal in $J$}

    Suppose now that $ p\in J$ and $ p$ is not maximal in $ J$. Because $p$ is minimal in $ \H{x\vee y}$, $ p$ is also minimal in $J$ since $J$ is a labeled order ideal of $ \H{x\vee y}$. Since $ p$ is minimal in $J$ with $ \H{y} \cong J$, it follows from Proposition \ref{removeminprop} part (1) that $s \in D_{L}(y) $. But note that because $y^{2}  =1$, it follows that $ D_{R}(y) = D_{L}(y)$. Thus, $ s \in D_{R}(y)$. But by Proposition \ref{removemaxprop} part (1), this implies that there exists some vertex $ z \in J \cong \H{y}$ such that $ z$ is maximal in $J$ and $ z$ has $S$-label $s$. Furthermore, $ z\neq p$ since $ p\in J$ is minimal but not maximal in $J$. Note that since $ z$ and $ q$ have the same $ S$-label in $ \H{x\vee y}$, it follows from Proposition \ref{mustbecompprop} that $ z \leq_{x\vee y} q$ or $ q \leq_{x\vee y} z$. But since $ q\in \H{x\vee y}$ is maximal, we conclude that $ z \leq_{x\vee y} q$. Since $ q\in I$ and $I$ is a labeled order ideal of $ \H{x\vee y}$, we conclude that $ z\in I$. Note that if $ z = q$, then one has $ \H{x\vee y} \setminus \{ p,q \} = (I \setminus \{  p,q \}) \cup (J \setminus \{  p,z \} ) $. If $ z\neq q$, then $ z <_{x\vee y} q$, and because $ z$ is maximal in $J$, we have that $ q\notin J$. Furthermore, because $ q\in I$, we still have $ \H{x\vee y} \setminus \{ p,q \} = (I \setminus \{ p,q \} ) \cup (J \setminus \{ p,z \}) $. Thus, we have shown that regardless of whether $ z=q$ or $ z\neq q$, we have $ \H{x\vee y} \setminus \{ p,q \} = (I \setminus \{ p,q \} ) \cup (J \setminus \{ p,z \}) $. Note that $I \setminus \{ p,q \} $ and $ J \setminus \{ p,z \}$ are labeled order ideals of $ \H{x\vee y} \setminus \{ p,q \}$. One can use Proposition \ref{removemaxprop} part (4) and Proposition \ref{removeminprop} part (4) to conclude that $ \H{s(x\vee y)s} \cong \H{x\vee y}\setminus \{ p,q \}$, $ \H{sxs} \cong I \setminus \{ p,q \}$, and $ \H{sys} \cong J \setminus \{  p,q \}$. Hence, one can Proposition \ref{joinidealprop} to deduce that $ s(x\vee y)s = (sxs)\vee (sys)$. But note that $ s(x\vee y)s$ is a fully commutative element where $ sxs \weakr s(x\vee y)s$, $ sys\weakr s(x\vee y)s$, $ (sxs)^{2} = 1$, $ (sys)^{2}=1$, and $ \ell(sxs)+\ell(sys) = \ell(x) - 2 + \ell(y) -2 = n+1 - 4 = n-3 < n+1$. Hence, we can apply induction to conclude that $ (s(x\vee y)s)^{2} = 1$, and thus $ (x\vee y)^{2} = 1$.

    The proof of the three cases above along with the proof of the case when $ \H{x\vee y}$ is composed solely of minimal elements exhausts all possible cases. Hence, we deduce that $ (x\vee y)^{2} = 1$ in general, which concludes that proof of the proposition.

\end{proof}

\end{document}
